% ==================================================================
%  Sujet 1 : Resine positive -- tests de dose et d'adhesion
%  Seance du 5 octobre 2026, DLL-MED -- assistante : S. Drljaca
% ==================================================================
\section{Résine positive -- tests de dose et d'adhésion}

\subsection{Objectif}
Ce sujet consiste en la déposition de la résine positive AZ1512 sur quatre wafers de silicium identiques, afin d’étudier l’influence de la vitesse de rotation sur l’épaisseur de la résine, mesurée à l’aide de l’instrument Filmetrics. On réalise également le transfert des motifs par photolithographie, à l’aide d’un photomasque et d’un rayonnement UV. Enfin, on caractérise l’hydrophobicité de la surface de silicium avec et sans un promoteur d'adhérence, en observant les angles de contact et réalisant des tests d’adhérence de la résine.


\subsection{Matériels}
\begin{table}[htbp]
  \centering\small
  \caption{Matériel utilisé pendant la séance.}
  \label{tab:s1-materiel}
  \begin{tabularx}{\linewidth}{@{}lX@{}}
    \toprule
    \textbf{Catégorie} & \textbf{Utilisé pendant la séance} \\
    \midrule
    Substrats   & 4 wafer de Si de 4 Pouces , échantillons de Si et film de Téflon pour la partie adhésion \\
    Produits    & Résine positive AZ1512, HMDS, développeur AZ 726 MIF, acétone, isopropanol (IPA), eau DI, azote \ce{N2} \\
    Équipements & Plaque chauffante, Spinner Laurell 400W, lampe UV P.250 en contact dur, Filmetrics F20, montage d'angle de contact (logiciel S-EYE 2.0), microscope optique Nikon, hotte \\
    \bottomrule
  \end{tabularx}
\end{table}

\subsection{Méthodes}

\paragraph{Partie A -- test de dose.}
\begin{enumerate}
  \item Déshydratation sur plaque chauffante (\qty{150}{\celsius}, \qty{5}{\minute}), puis HMDS à la tournette (\qty{3000}{\tour\per\minute}, \qty{30}{\second}), si nécessaire.
  \item Étalement de l'AZ1512 : \qty{700}{\tour\per\minute} pendant \qty{10}{\second}, puis \numlist{2000;3000;4000;5000}~\unit{\tour\per\minute} selon la plaquette.
  \item Recuit doux à \qty{100}{\celsius} pendant \qty{60}{\second}, puis refroidissement : la résine reprend l'humidité de l'air nécessaire à la photoréaction.
  \item Exposition en contact dur sous la lampe P.250 ($I = \qty{3,11}{\milli\watt\per\centi\metre\squared}$) de 8 zones périphériques. Temps : \qtyrange{1}{8}{\second} par pas de \qty{1}{\second} pour 2000 et \qty{3000}{\tour\per\minute} ; \qtyrange{0,5}{3,5}{\second} par pas de \qty{0,5}{\second}, puis \qty{5}{\second}, pour 4000 et \qty{5000}{\tour\per\minute}.
  \item Développement dans l'AZ 726 MIF pendant \qty{40}{\second} avec agitation, rinçage à l'eau DI, séchage à l'azote.
  \item Mesure au Filmetrics F20 : épaisseur initiale \ez{} en 5 points centraux non exposés (O, A, B, C, D) et épaisseur restante $e$ dans les 8 zones exposées.
\end{enumerate}

\paragraph{Mesure d'épaisseur par réflectométrie spectrale.}
La lumière blanche arrive à incidence normale et se réfléchit sur l'interface air/résine et sur l'interface résine/silicium. Les deux ondes interfèrent : la réflectance est maximale quand $2\,n\,e = m\,\lambda$ ($m$ entier), ce qui produit des franges dans le spectre. Le logiciel ajuste $e$ et l'indice $n(\lambda)$ (modèle de Cauchy) jusqu'à ce que le spectre calculé recouvre le spectre mesuré. Le GOF (\emph{goodness of fit}, de 0 à 1) indique la qualité de cet ajustement.

\paragraph{Exploitation.}
La dose reçue par chaque zone vaut $D = I\,t$ (\unit{\milli\watt\per\centi\metre\squared} $\times$ \unit{\second} $=$ \unit{\milli\joule\per\centi\metre\squared}). On trace l'épaisseur normalisée $e/\ez$ en fonction de $\log D$ et on ajuste la zone de transition par une droite. Son extrapolation à $e/\ez = 1$ donne \Dc, la dose où la dissolution commence ; son extrapolation à $e/\ez = 0$ donne \Dz, la dose de clarification, qui définit la sensibilité. Nous suivons la notation des slides d'introduction \cite{intro} ; le protocole \cite{prot1} appelle \enquote{$D_{100}$} la dose de clarification, c'est-à-dire notre \Dz. Le contraste est l'opposé de la pente de la droite :
\begin{equation}
  \gamma = \frac{1}{\log_{10}\left(D_0 / D_{100}\right)} .
  \label{eq:contraste}
\end{equation}

\paragraph{Partie B -- traitement de surface et test d'adhésion.}
On compare un échantillon traité au HMDS et un échantillon non traité.
\begin{enumerate}
  \item Angle de contact sur l'échantillon tel que reçu, puis nettoyage à l'acétone et à l'IPA (sans laisser sécher l'acétone) et séchage à l'azote.
  \item Déshydratation (\qty{150}{\celsius}, \qtyrange{5}{10}{\minute}), refroidissement, angle de contact.
  \item HMDS à \qty{4000}{\tour\per\minute} pendant \qty{30}{\second}, recuit à \qty{120}{\celsius} pendant \qty{10}{\minute}, angle de contact.
  \item AZ1512 (\qty{700}{\tour\per\minute} pendant \qty{10}{\second}, puis \qty{3000}{\tour\per\minute} pendant \qty{60}{\second}), recuit doux à \qty{100}{\celsius} pendant \qty{90}{\second}, exposition de \qty{9}{\second} (P.250, contact dur), développement de \qty{2}{\minute}, rinçage, séchage, inspection au microscope.
  \item Test d'adhésion : eau à \qty{100}{\celsius} pendant \qty{5}{\minute}, développement de \qty{30}{\second}, séchage, inspection. Angle de contact du Téflon comme référence hydrophobe.
\end{enumerate}

\paragraph{Mesure de l'angle de contact.}
Une goutte d'eau DI est déposée à la seringue et photographiée de profil (S-EYE 2.0). L'angle est mesuré dans le liquide, entre la surface et la tangente au profil au point de contact, et moyenné sur au moins trois gouttes.

\subsection{Résultats}


\begin{table}[H]
  \centering\footnotesize
  \setlength{\tabcolsep}{4pt}
  \caption{Paramètres d'enduction, épaisseurs initiales mesurées au Filmetrics (GOF entre parenthèses) et grandeurs extraites des courbes de contraste pour les wafers 1 à 4 (résine AZ1512).}
  \label{tab:topic1-epaisseurs}
  \begin{tabularx}{\linewidth}{@{}>{\raggedright\arraybackslash}X c c c c c@{}}
    \toprule
    \textbf{Grandeur} & \textbf{Unité} & \textbf{Wafer 1} & \textbf{Wafer 2} & \textbf{Wafer 3} & \textbf{Wafer 4} \\
    & & (Olivia) & (Titouan DLB) & (Lilian) & (Titouan B) \\
    \midrule
    Résine & --- & AZ1512 & AZ1512 & AZ1512 & AZ1512 \\
    Puissance de la lampe P250 & \unit{\milli\watt\per\centi\metre\squared} & 3,11 & 3,11 & 3,11 & 3,11 \\
    Vitesse de rotation finale & rpm & 2000 & 3000 & 4000 & 5000 \\
    Temps de rotation (2\textsuperscript{e} étape) & \unit{\second} & 90 & 90 & 90 & 90 \\
    \midrule
    Épaisseur zone O (GOF) & \unit{\micro\metre} & 1,8128 (0,938) & 1,4245 (0,998) & 1,2996 (0,999) & 1,2994 (0,999) \\
    Épaisseur zone A (GOF) & \unit{\micro\metre} & 1,7980 (0,999) & 1,4227 (0,999) & 1,2966 (0,999) & 1,2965 (0,999) \\
    Épaisseur zone B (GOF) & \unit{\micro\metre} & 1,8184 (0,996) & 1,4195 (0,999) & 1,1529 (0,999) & 1,2964 (0,999) \\
    Épaisseur zone C (GOF) & \unit{\micro\metre} & 1,8708 (0,809) & 1,4304 (0,999) & 1,1550 (0,999) & 1,2915 (0,999) \\
    Épaisseur zone D (GOF) & \unit{\micro\metre} & 1,7966 (0,999) & 1,4240 (0,999) & 1,1523 (0,999) & 1,2941 (0,999) \\
    Épaisseur moyenne $E_i$ & \unit{\micro\metre} & 1,819 & 1,424 & 1,211 & 1,296 \\
    Écart-type $\sigma$ & \unit{\micro\metre} & 0,030 & 0,004 & 0,079 & 0,003 \\
    \midrule
    Dose de début de dissolution $D_0$ & \unit{\milli\joule\per\centi\metre\squared} & 3,2 & 2,9 & --- & --- \\
    Dose de dégagement $D_{100}$ & \unit{\milli\joule\per\centi\metre\squared} & 19 & 13 & 1,6 à 3,1 & $<$ 1,6 \\
    Contraste $\gamma$ & --- & 1,3 & 1,5 & n.m. & n.m. \\
    \bottomrule
  \end{tabularx}

  \vspace{2pt}
  {\scriptsize n.m. : non mesurable (aucun point dans la zone de transition).}
\end{table}

\begin{table}[H]
  \centering\footnotesize
  \setlength{\tabcolsep}{4pt}
  \caption{Temps d'insolation, dose, épaisseur finale après développement et épaisseur normalisée pour les 8 zones exposées de chaque wafer ($I = \qty{3,11}{\milli\watt\per\centi\metre\squared}$).}
  \label{tab:topic1-dose}
  \begin{tabularx}{\linewidth}{@{}>{\raggedright\arraybackslash}X c c c c c c c c c@{}}
    \toprule
    \textbf{Grandeur} & \textbf{Unité} & \textbf{1} & \textbf{2} & \textbf{3} & \textbf{4} & \textbf{5} & \textbf{6} & \textbf{7} & \textbf{8} \\
    \midrule
    \multicolumn{10}{@{}l}{\textbf{Wafer 1} (2000 rpm, $E_i = \qty{1,819}{\micro\metre}$)} \\
    Temps d'insolation $t_{\text{exp}}$ & \unit{\second} & 2 & 3 & 4 & 5 & 6 & 7 & 8 & 9 \\
    Dose $D = I \times t_{\text{exp}}$ & \unit{\milli\joule\per\centi\metre\squared} & 6,22 & 9,33 & 12,44 & 15,55 & 18,66 & 21,77 & 24,88 & 27,99 \\
    Épaisseur finale $E_f$ & \unit{\micro\metre} & 1,126 & 0,778 & 0,440 & 0,173 & 0,044 & 0,005 & 0,004 & 0,002 \\
    Épaisseur normalisée $E_f/E_i$ & --- & 0,619 & 0,428 & 0,242 & 0,095 & 0,024 & 0,003 & 0,002 & 0,001 \\
    \midrule
    \multicolumn{10}{@{}l}{\textbf{Wafer 2} (3000 rpm, $E_i = \qty{1,424}{\micro\metre}$)} \\
    Temps d'insolation $t_{\text{exp}}$ & \unit{\second} & 1 & 2 & 3 & 4 & 5 & 6 & 7 & 8 \\
    Dose $D = I \times t_{\text{exp}}$ & \unit{\milli\joule\per\centi\metre\squared} & 3,11 & 6,22 & 9,33 & 12,44 & 15,55 & 18,66 & 21,77 & 24,88 \\
    Épaisseur finale $E_f$ & \unit{\micro\metre} & 1,298 & 0,808 & 0,347 & 0,016 & 0,026 & 0,011 & 0,008 & 0,006 \\
    Épaisseur normalisée $E_f/E_i$ & --- & 0,912 & 0,567 & 0,243 & 0,011 & 0,018 & 0,008 & 0,005 & 0,004 \\
    \midrule
    \multicolumn{10}{@{}l}{\textbf{Wafer 3} (4000 rpm, $E_i = \qty{1,211}{\micro\metre}$)} \\
    Temps d'insolation $t_{\text{exp}}$ & \unit{\second} & 0,5 & 1 & 1,5 & 2 & 2,5 & 3 & 3,5 & 5 \\
    Dose $D = I \times t_{\text{exp}}$ & \unit{\milli\joule\per\centi\metre\squared} & 1,56 & 3,11 & 4,67 & 6,22 & 7,78 & 9,33 & 10,89 & 15,55 \\
    Épaisseur finale $E_f$ & \unit{\micro\metre} & 0,725 & 0,028 & 0,004 & 0,006 & 0,003 & 0,000 & 0,004 & 0,006 \\
    Épaisseur normalisée $E_f/E_i$ & --- & 0,598 & 0,023 & 0,003 & 0,005 & 0,003 & 0,000 & 0,003 & 0,005 \\
    \midrule
    \multicolumn{10}{@{}l}{\textbf{Wafer 4} (5000 rpm, $E_i = \qty{1,296}{\micro\metre}$)} \\
    Temps d'insolation $t_{\text{exp}}$ & \unit{\second} & 0,5 & 1 & 1,5 & 2 & 2,5 & 3 & 3,5 & 5 \\
    Dose $D = I \times t_{\text{exp}}$ & \unit{\milli\joule\per\centi\metre\squared} & 1,56 & 3,11 & 4,67 & 6,22 & 7,78 & 9,33 & 10,89 & 15,55 \\
    Épaisseur finale $E_f$ & \unit{\micro\metre} & 0,008 & 0,007 & 0,003 & 0,005 & 0,013 & 0,006 & 0,007 & 0,006 \\
    Épaisseur normalisée $E_f/E_i$ & --- & 0,006 & 0,005 & 0,002 & 0,004 & 0,010 & 0,005 & 0,006 & 0,004 \\
    \bottomrule
  \end{tabularx}
\end{table}
\pgfplotsset{
  topic1/.style={
    width=0.5\linewidth, height=4.2cm,
    tick label style={font=\scriptsize}, label style={font=\footnotesize},
    title style={font=\footnotesize, yshift=-3pt},
    /pgf/number format/use comma,
    grid=major, grid style={gray!20},
    tick align=outside, tick pos=left,
  },
  contraste/.style={
    topic1, xmin=1, xmax=30, ymin=0, ymax=1.05,
    xtick={1,2,5,10,20,30}, xticklabels={1,2,5,10,20,30},
    log ticks with fixed point,
    ytick={0,0.2,...,1}, yticklabel style={/pgf/number format/fixed, /pgf/number format/precision=1},
  },
}

% ===============================================================
% Figure 1 : courbes de contraste des 4 wafers
% ===============================================================
\begin{figure}[H]
  \centering
  \begin{tikzpicture}
  \begin{groupplot}[
    group style={group size=2 by 2, horizontal sep=1.2cm, vertical sep=1.4cm,
                 xlabels at=edge bottom, ylabels at=edge left},
    contraste, xmode=log,
    xlabel={Dose (\unit{\milli\joule\per\centi\metre\squared})},
    ylabel={$E_f/E_i$},
  ]
    \nextgroupplot[title={Wafer 1 (2000 rpm, $E_i = \qty{1,82}{\micro\metre}$)}]
    \addplot[gray, dashed, thin] coordinates {(3.217,1) (19.157,0)};
    \draw[gray, thin] (axis cs:19.157,0) -- (axis cs:19.157,0.28);
    \node[anchor=north east, font=\scriptsize, align=right] at (rel axis cs:0.97,0.95) {$D_{100} \approx \qty{19}{\milli\joule\per\centi\metre\squared}$\\ $\gamma \approx 1{,}3$};
    \addplot[black, mark=*, mark size=1.3pt, thin] coordinates {(6.2200,0.6187) (9.3300,0.4277) (12.4400,0.2416) (15.5500,0.0950) (18.6600,0.0242) (21.7700,0.0029) (24.8800,0.0021) (27.9900,0.0009)};
    \nextgroupplot[title={Wafer 2 (3000 rpm, $E_i = \qty{1,42}{\micro\metre}$)}]
    \addplot[gray, dashed, thin] coordinates {(2.866,1) (13.451,0)};
    \draw[gray, thin] (axis cs:13.451,0) -- (axis cs:13.451,0.28);
    \node[anchor=north east, font=\scriptsize, align=right] at (rel axis cs:0.97,0.95) {$D_{100} \approx \qty{13}{\milli\joule\per\centi\metre\squared}$\\ $\gamma \approx 1{,}5$};
    \addplot[black, mark=*, mark size=1.3pt, thin] coordinates {(3.1100,0.9116) (6.2200,0.5674) (9.3300,0.2434) (12.4400,0.0109) (15.5500,0.0183) (18.6600,0.0077) (21.7700,0.0054) (24.8800,0.0043)};
    \nextgroupplot[title={Wafer 3 (4000 rpm, $E_i = \qty{1,21}{\micro\metre}$)}]
    \fill[gray!18] (axis cs:1.555,0) rectangle (axis cs:3.110,1.05);
    \node[anchor=north west, font=\scriptsize] at (axis cs:1.555,1.02) {$D_{100}$ entre 1,6 et 3,1};
    \addplot[black, mark=*, mark size=1.3pt, thin] coordinates {(1.5550,0.5985) (3.1100,0.0233) (4.6650,0.0031) (6.2200,0.0053) (7.7750,0.0026) (9.3300,0.0000) (10.8850,0.0033) (15.5500,0.0049)};
    \nextgroupplot[title={Wafer 4 (5000 rpm, $E_i = \qty{1,30}{\micro\metre}$)}]
    \fill[gray!18] (axis cs:1.000,0) rectangle (axis cs:1.555,1.05);
    \node[anchor=north west, font=\scriptsize] at (axis cs:1.000,1.02) {$D_{100} < 1{,}6$};
    \addplot[black, mark=*, mark size=1.3pt, thin] coordinates {(1.5550,0.0064) (3.1100,0.0053) (4.6650,0.0024) (6.2200,0.0042) (7.7750,0.0098) (9.3300,0.0048) (10.8850,0.0056) (15.5500,0.0043)};
  \end{groupplot}
  \end{tikzpicture}
  \caption{Courbes de contraste de l'AZ1512 pour les 4 wafers : épaisseur normalisée $E_f/E_i$ en fonction de la dose (échelle logarithmique). En pointillés, l'ajustement de la zone linéaire (wafers 1 et 2) ; en gris, l'encadrement de $D_{100}$ (wafers 3 et 4).}
  \label{fig:topic1-contraste}
\end{figure}

% ===============================================================
% Figure 2 : courbe d'étalement et D100 en fonction de l'épaisseur
% ===============================================================
\begin{figure}[H]
  \centering
  \begin{tikzpicture}
  \begin{groupplot}[
    group style={group size=2 by 1, horizontal sep=1.6cm},
    topic1,
  ]
    % (a) épaisseur en fonction de la vitesse
    \nextgroupplot[title={(a) Courbe d'étalement},
      xlabel={Vitesse de rotation (rpm)}, ylabel={$E_i$ (\unit{\micro\metre})},
      xmin=1500, xmax=5500, ymin=1, ymax=2,
      xtick={2000,3000,4000,5000},
      xticklabel style={/pgf/number format/1000 sep={}},
      yticklabel style={/pgf/number format/fixed, /pgf/number format/precision=1}]
    \addplot[gray, dashed, thin, domain=1800:5200, samples=60] {78.66/sqrt(x)};
    \addplot[black, only marks, mark=*, mark size=1.4pt,
             error bars/.cd, y dir=both, y explicit]
      coordinates {(2000,1.81932) +- (0,0.03027) (3000,1.42422) +- (0,0.00397) (4000,1.21128) +- (0,0.07927) (5000,1.29558) +- (0,0.00296)};
    \node[anchor=north east, font=\scriptsize, gray, align=right] at (rel axis cs:0.97,0.95) {pointillés : loi en $1/\sqrt{\omega}$\\ ajustée sur les wafers 1 à 3};

    % (b) D100 en fonction de l'épaisseur initiale
    \nextgroupplot[title={(b) Sensibilité $D_{100}$},
      xlabel={$E_i$ (\unit{\micro\metre})}, ylabel={$D_{100}$ (\unit{\milli\joule\per\centi\metre\squared})},
      xmin=1.1, xmax=1.9, ymin=0, ymax=22,
      xtick={1.2,1.4,1.6,1.8},
      xticklabel style={/pgf/number format/fixed, /pgf/number format/precision=1}]
    \addplot[gray, dashed, thin, domain=1.1:1.9] {13.451+(x-1.42422)*14.442};
    \addplot[black, only marks, mark=*, mark size=1.4pt] coordinates {(1.81932,19.16) (1.42422,13.45)};
    \addplot[black, only marks, mark=*, mark size=1.4pt, error bars/.cd, y dir=both, y explicit]
      coordinates {(1.21128,2.333) +- (0,0.777)};
    \addplot[black, only marks, mark=triangle*, mark size=2pt, every mark/.append style={rotate=180}]
      coordinates {(1.29558,1.555)};
    \node[anchor=west, font=\scriptsize] at (axis cs:1.81932,19.16) {W1};
    \node[anchor=north west, font=\scriptsize] at (axis cs:1.42422,13.45) {W2};
    \node[anchor=south, font=\scriptsize] at (axis cs:1.21128,3.110) {W3};
    \node[anchor=west, font=\scriptsize] at (axis cs:1.29558,1.555) {W4 ($<$)};
  \end{groupplot}
  \end{tikzpicture}
  \caption{(a) Épaisseur initiale moyenne en fonction de la vitesse de rotation (barres : écart-type). (b) Dose de dégagement $D_{100}$ en fonction de l'épaisseur initiale ; pour les wafers 3 et 4, seul un encadrement est possible. En pointillés, la tendance des wafers 1 et 2.}
  \label{fig:topic1-etalement-d100}
\end{figure}

% ===============================================================
% Figure 3 : courbes attendues et mesurées des wafers 3 et 4
% ===============================================================
\begin{figure}[H]
  \centering
  \begin{tikzpicture}
  \begin{axis}[contraste, xmode=log, width=0.62\linewidth, height=4.4cm,
    xlabel={Dose (\unit{\milli\joule\per\centi\metre\squared})}, ylabel={$E_f/E_i$},
    legend style={font=\scriptsize, draw=none, fill=none, at={(0.99,0.98)}, anchor=north east},
    legend cell align=left]
    \addplot[black, dashed, thin, domain=1:30, samples=120] {max(0,min(1,1.390*log10(10.376/x)))};
    \addlegendentry{Wafer 3 attendu}
    \addplot[gray, dashed, thin, domain=1:30, samples=120] {max(0,min(1,1.390*log10(11.594/x)))};
    \addlegendentry{Wafer 4 attendu}
    \addplot[black, mark=*, mark size=1.3pt, thin] coordinates {(1.5550,0.5985) (3.1100,0.0233) (4.6650,0.0031) (6.2200,0.0053) (7.7750,0.0026) (9.3300,0.0000) (10.8850,0.0033) (15.5500,0.0049)};
    \addlegendentry{Wafer 3 mesuré}
    \addplot[gray, mark=square*, mark size=1.2pt, thin] coordinates {(1.5550,0.0064) (3.1100,0.0053) (4.6650,0.0024) (6.2200,0.0042) (7.7750,0.0098) (9.3300,0.0048) (10.8850,0.0056) (15.5500,0.0043)};
    \addlegendentry{Wafer 4 mesuré}
  \end{axis}
  \end{tikzpicture}
  \caption{Courbes de contraste attendues (pointillés, à partir de la tendance de $D_{100}$ et du $\gamma$ moyen des wafers 1 et 2) et mesurées pour les wafers 3 et 4.}
  \label{fig:topic1-attendu}
\end{figure}
\subsection{Analyse}

\paragraph{Épaisseur et vitesse de rotation}

Entre 2000 et 4000~rpm, l'épaisseur baisse de 1,82 à 1,21~$\mu$m, en suivant à peu près la loi attendue en $1/\sqrt{\omega}$. Le wafer~4 (1,30~$\mu$m à 5000~rpm) est plus épais que le wafer~3, ce qui n'est pas physique : nous pensons à une inversion de wafers ou à une vitesse mal réglée pendant le traitement en parallèle. Les films sont uniformes (écart-type de 3 à 30~nm), sauf le wafer~3 (79~nm) dont les zones O et A sont plus épaisses.

\paragraph{Comparaison des courbes de contraste}

Pour les wafers 1 et 2, $D_{100}$ augmente avec l'épaisseur (19 contre 13~mJ/cm$^2$) : la résine absorbe les UV, donc le bas d'un film épais reçoit moins de dose. Le contraste reste proche ($\gamma = 1{,}3$ et $1{,}5$), car il dépend surtout de la résine et du développeur.

Les courbes plates des wafers 3 et 4 sont le point principal de la séance. Le protocole donne des temps d'exposition à adapter selon les résultats des premiers wafers, mais nous avons traité les quatre wafers en parallèle avec l'assistante. Les wafers 3 et 4 ont donc reçu directement les temps suggérés (0,5 à 5~s) : presque toutes leurs zones ont dépassé $D_{100}$, la transition n'a pas été mesurée, et on n'obtient qu'un encadrement de $D_{100}$ et aucun $\gamma$. Ces $D_{100}$ sont aussi bien plus bas que la tendance des wafers 1 et 2 (environ 10 à 12~mJ/cm$^2$), donc l'imprécision des temps très courts sur la lampe P250 ou des différences de développement ont sans doute joué aussi. Il aurait fallu développer et mesurer chaque wafer avant d'exposer le suivant, puis recentrer les temps autour de la transition. Ici, on a donc une courbe complètement plate et rien à quantifier.

\subsection{Conclusion}
L'épaisseur de l'AZ1512 passe de \qty{1,82}{\micro\metre} à \qty{1,21}{\micro\metre} entre 2000 et \qty{4000}{\tour\per\minute}, conformément à la loi en $\omega^{-1/2}$ ; la plaquette à \qty{5000}{\tour\per\minute} sort de cette tendance. Pour les deux films les plus épais, la sensibilité est proportionnelle à l'épaisseur (\qtylist{13,5;17,4}{\milli\joule\per\centi\metre\squared}) et le contraste vaut 0,9 à 1,5. Les films minces se sont développés à des doses trop faibles pour notre grille de temps : il faut adapter les temps après une première zone test. 
% ==================================================================
%  Sujet 1 : Resine positive -- tests de dose et d'adhesion
%  Seance du 5 octobre 2026, DLL-MED -- assistante : S. Drljaca
% ==================================================================
\section{Résine positive -- tests de dose et d'adhésion}

\subsection{Objectif}
Ce sujet consiste à déposer la résine positive AZ1512 sur quatre wafers de silicium identiques, afin d'étudier l'influence de la vitesse de rotation sur l'épaisseur de la résine, mesurée à l'aide du Filmetrics. On réalise également le transfert de motifs par photolithographie, à l'aide d'un photomasque et d'un rayonnement UV, pour tracer la courbe de contraste de la résine. Enfin, on caractérise l'hydrophobicité de la surface de silicium avec et sans promoteur d'adhérence, en mesurant les angles de contact et en réalisant des tests d'adhérence de la résine.

\subsection{Matériels}
\begin{table}[H]
  \centering\small
  \caption{Matériel utilisé pendant la séance.}
  \label{tab:s1-materiel}
  \begin{tabularx}{\linewidth}{@{}lX@{}}
    \toprule
    \textbf{Catégorie} & \textbf{Utilisé pendant la séance} \\
    \midrule
    Substrats   & 4 wafers de Si de 4~pouces, échantillons de Si et film de Téflon pour la partie adhésion \\
    Produits    & Résine positive AZ1512, HMDS, développeur AZ 726 MIF, acétone, isopropanol (IPA), eau DI, azote \ce{N2} \\
    Équipements & Plaque chauffante, spinner Laurell 400W, lampe UV P.250 en contact dur, Filmetrics F20, montage d'angle de contact (logiciel S-EYE 2.0), microscope optique Nikon, hotte \\
    \bottomrule
  \end{tabularx}
\end{table}

\subsection{Méthodes}

\paragraph{Partie A -- test de dose.}
\begin{enumerate}
  \item Déshydratation (\qty{150}{\celsius}, \qty{5}{\minute}) et HMDS à la tournette (\qty{3000}{\tour\per\minute}, \qty{30}{\second}), réalisés avant la séance.
  \item Étalement de l'AZ1512 : \qty{700}{\tour\per\minute} pendant \qty{10}{\second}, puis \numlist{2000;3000;4000;5000}~\unit{\tour\per\minute} selon le wafer, pendant \qty{90}{\second} (feuille de mesure ; le protocole indique \qty{60}{\second}).
  \item Recuit doux à \qty{100}{\celsius} pendant \qty{60}{\second}, puis refroidissement : la résine reprend l'humidité de l'air nécessaire à la photoréaction.
  \item Exposition en contact dur sous la lampe P.250 ($I = \qty{3,11}{\milli\watt\per\centi\metre\squared}$) de 8 zones périphériques. Temps : \qtyrange{2}{9}{\second} (\qty{2000}{\tour\per\minute}) et \qtyrange{1}{8}{\second} (\qty{3000}{\tour\per\minute}) par pas de \qty{1}{\second} ; \qtyrange{0,5}{3,5}{\second} par pas de \qty{0,5}{\second} puis \qty{5}{\second} (4000 et \qty{5000}{\tour\per\minute}).
  \item Développement dans l'AZ 726 MIF pendant \qty{40}{\second} avec agitation, rinçage à l'eau DI, séchage à l'azote.
  \item Mesure au Filmetrics F20 : épaisseur initiale \ez{} en 5 points centraux non exposés (O, A, B, C, D) et épaisseur restante $e$ dans les 8 zones exposées.
\end{enumerate}

\paragraph{Mesure d'épaisseur par réflectométrie spectrale.}
La lumière blanche arrive en incidence normale et se réfléchit sur l'interface air/résine et sur l'interface résine/silicium. Les deux ondes interfèrent : la réflectance est maximale quand $2\,n\,e = m\,\lambda$ ($m$ entier), ce qui produit des franges dans le spectre. Le logiciel ajuste $e$ et l'indice $n(\lambda)$ (modèle de Cauchy) jusqu'à ce que le spectre calculé recouvre le spectre mesuré. Le GOF (\emph{goodness of fit}, de 0 à 1) indique la qualité de cet ajustement.

\paragraph{Exploitation.}
La dose reçue par chaque zone vaut $D = I\,t$ (\unit{\milli\watt\per\centi\metre\squared} $\times$ \unit{\second} $=$ \unit{\milli\joule\per\centi\metre\squared}). On trace l'épaisseur normalisée $e/\ez$ en fonction de $\log D$ et on ajuste la zone de transition par une droite. Son extrapolation à $e/\ez = 1$ donne \Dc, la dose où la dissolution commence ; son extrapolation à $e/\ez = 0$ donne \Dz, la dose de clarification, qui définit la sensibilité. Nous suivons la notation des slides d'introduction \cite{intro} ; le protocole \cite{prot1} appelle \enquote{$D_{100}$} la dose de clarification, c'est-à-dire notre \Dz. Le contraste est l'opposé de la pente de la droite :
\begin{equation}
  \gamma = \frac{1}{\log_{10}\left(\Dz / \Dc\right)} .
  \label{eq:contraste}
\end{equation}

\paragraph{Partie B -- traitement de surface et test d'adhésion.}
On compare deux échantillons de Si : l'un suit toutes les étapes ci-dessous, l'autre ne reçoit pas le HMDS (étape~3).
\begin{enumerate}
  \item Angle de contact sur l'échantillon tel que reçu, puis nettoyage à l'acétone et à l'IPA (sans laisser sécher l'acétone) et séchage à l'azote.
  \item Déshydratation (\qty{150}{\celsius}, \qtyrange{5}{10}{\minute}), refroidissement, angle de contact.
  \item HMDS à \qty{4000}{\tour\per\minute} pendant \qty{30}{\second}, recuit à \qty{120}{\celsius} pendant \qty{10}{\minute}, angle de contact.
  \item AZ1512 (\qty{700}{\tour\per\minute} pendant \qty{10}{\second}, puis \qty{3000}{\tour\per\minute} pendant \qty{60}{\second}), recuit doux à \qty{100}{\celsius} pendant \qty{90}{\second}, exposition de \qty{9}{\second} (P.250, contact dur), développement de \qty{2}{\minute}, rinçage, séchage, inspection au microscope.
  \item Test d'adhésion : eau à \qty{100}{\celsius} pendant \qty{5}{\minute}, développement de \qty{30}{\second}, séchage, inspection. Angle de contact du Téflon comme référence hydrophobe.
\end{enumerate}

\paragraph{Mesure de l'angle de contact.}
Une goutte d'eau DI est déposée à la seringue et photographiée de profil (S-EYE 2.0). L'angle est mesuré dans le liquide, entre la surface et la tangente au profil au point de contact, et moyenné sur au moins trois gouttes.

\subsection{Résultats}

Les mesures sont regroupées dans les tableaux~\ref{tab:topic1-epaisseurs} et~\ref{tab:topic1-dose}, et les courbes correspondantes aux figures~\ref{fig:topic1-contraste} et~\ref{fig:topic1-etalement-d100}.

\begin{table}[H]
  \centering\footnotesize
  \setlength{\tabcolsep}{4pt}
  \caption{Paramètres d'enduction, épaisseurs initiales mesurées au Filmetrics (GOF entre parenthèses) et grandeurs extraites des courbes de contraste pour les wafers 1 à 4 (résine AZ1512).}
  \label{tab:topic1-epaisseurs}
  \begin{tabularx}{\linewidth}{@{}>{\raggedright\arraybackslash}X c c c c c@{}}
    \toprule
    \textbf{Grandeur} & \textbf{Unité} & \textbf{Wafer 1} & \textbf{Wafer 2} & \textbf{Wafer 3} & \textbf{Wafer 4} \\
    & & (Olivia) & (Titouan DLB) & (Lilian) & (Titouan B) \\
    \midrule
    Résine & --- & AZ1512 & AZ1512 & AZ1512 & AZ1512 \\
    Intensité de la lampe P.250 & \unit{\milli\watt\per\centi\metre\squared} & 3,11 & 3,11 & 3,11 & 3,11 \\
    Vitesse de rotation finale & \unit{\tour\per\minute} & 2000 & 3000 & 4000 & 5000 \\
    Temps de rotation (2\textsuperscript{e} étape) & \unit{\second} & 90 & 90 & 90 & 90 \\
    \midrule
    Épaisseur zone O (GOF) & \unit{\micro\metre} & 1,8128 (0,938) & 1,4245 (0,998) & 1,2996 (0,999) & 1,2994 (0,999) \\
    Épaisseur zone A (GOF) & \unit{\micro\metre} & 1,7980 (0,999) & 1,4227 (0,999) & 1,2966 (0,999) & 1,2965 (0,999) \\
    Épaisseur zone B (GOF) & \unit{\micro\metre} & 1,8184 (0,996) & 1,4195 (0,999) & 1,1529 (0,999) & 1,2964 (0,999) \\
    Épaisseur zone C (GOF) & \unit{\micro\metre} & 1,8708 (0,809) & 1,4304 (0,999) & 1,1550 (0,999) & 1,2915 (0,999) \\
    Épaisseur zone D (GOF) & \unit{\micro\metre} & 1,7966 (0,999) & 1,4240 (0,999) & 1,1523 (0,999) & 1,2941 (0,999) \\
    Épaisseur initiale moyenne \ez & \unit{\micro\metre} & 1,819 & 1,424 & 1,211 & 1,296 \\
    Écart-type $\sigma$ & \unit{\micro\metre} & 0,030 & 0,004 & 0,079 & 0,003 \\
    \midrule
    Dose de début de dissolution \Dc & \unit{\milli\joule\per\centi\metre\squared} & 3,2 & 2,9 & --- & --- \\
    Dose de clarification \Dz & \unit{\milli\joule\per\centi\metre\squared} & 19 & 13 & 1,6 à 3,1 & $<$ 1,6 \\
    Contraste $\gamma$ & --- & 1,3 & 1,5 & n.m. & n.m. \\
    \bottomrule
  \end{tabularx}

  \vspace{2pt}
  {\scriptsize n.m. : non mesurable (aucun point dans la zone de transition).}
\end{table}

\begin{table}[H]
  \centering\footnotesize
  \setlength{\tabcolsep}{4pt}
  \caption{Temps d'insolation, dose, épaisseur restante après développement et épaisseur normalisée pour les 8 zones exposées de chaque wafer ($I = \qty{3,11}{\milli\watt\per\centi\metre\squared}$).}
  \label{tab:topic1-dose}
  \begin{tabularx}{\linewidth}{@{}>{\raggedright\arraybackslash}X c c c c c c c c c@{}}
    \toprule
    \textbf{Grandeur} & \textbf{Unité} & \textbf{1} & \textbf{2} & \textbf{3} & \textbf{4} & \textbf{5} & \textbf{6} & \textbf{7} & \textbf{8} \\
    \midrule
    \multicolumn{10}{@{}l}{\textbf{Wafer 1} (\qty{2000}{\tour\per\minute}, $\ez = \qty{1,819}{\micro\metre}$)} \\
    Temps d'insolation $t$ & \unit{\second} & 2 & 3 & 4 & 5 & 6 & 7 & 8 & 9 \\
    Dose $D = I\,t$ & \unit{\milli\joule\per\centi\metre\squared} & 6,22 & 9,33 & 12,44 & 15,55 & 18,66 & 21,77 & 24,88 & 27,99 \\
    Épaisseur restante $e$ & \unit{\micro\metre} & 1,126 & 0,778 & 0,440 & 0,173 & 0,044 & 0,005 & 0,004 & 0,002 \\
    Épaisseur normalisée $e/\ez$ & --- & 0,619 & 0,428 & 0,242 & 0,095 & 0,024 & 0,003 & 0,002 & 0,001 \\
    \midrule
    \multicolumn{10}{@{}l}{\textbf{Wafer 2} (\qty{3000}{\tour\per\minute}, $\ez = \qty{1,424}{\micro\metre}$)} \\
    Temps d'insolation $t$ & \unit{\second} & 1 & 2 & 3 & 4 & 5 & 6 & 7 & 8 \\
    Dose $D = I\,t$ & \unit{\milli\joule\per\centi\metre\squared} & 3,11 & 6,22 & 9,33 & 12,44 & 15,55 & 18,66 & 21,77 & 24,88 \\
    Épaisseur restante $e$ & \unit{\micro\metre} & 1,298 & 0,808 & 0,347 & 0,016 & 0,026 & 0,011 & 0,008 & 0,006 \\
    Épaisseur normalisée $e/\ez$ & --- & 0,912 & 0,567 & 0,243 & 0,011 & 0,018 & 0,008 & 0,005 & 0,004 \\
    \midrule
    \multicolumn{10}{@{}l}{\textbf{Wafer 3} (\qty{4000}{\tour\per\minute}, $\ez = \qty{1,211}{\micro\metre}$)} \\
    Temps d'insolation $t$ & \unit{\second} & 0,5 & 1 & 1,5 & 2 & 2,5 & 3 & 3,5 & 5 \\
    Dose $D = I\,t$ & \unit{\milli\joule\per\centi\metre\squared} & 1,56 & 3,11 & 4,67 & 6,22 & 7,78 & 9,33 & 10,89 & 15,55 \\
    Épaisseur restante $e$ & \unit{\micro\metre} & 0,725 & 0,028 & 0,004 & 0,006 & 0,003 & 0,000 & 0,004 & 0,006 \\
    Épaisseur normalisée $e/\ez$ & --- & 0,598 & 0,023 & 0,003 & 0,005 & 0,003 & 0,000 & 0,003 & 0,005 \\
    \midrule
    \multicolumn{10}{@{}l}{\textbf{Wafer 4} (\qty{5000}{\tour\per\minute}, $\ez = \qty{1,296}{\micro\metre}$)} \\
    Temps d'insolation $t$ & \unit{\second} & 0,5 & 1 & 1,5 & 2 & 2,5 & 3 & 3,5 & 5 \\
    Dose $D = I\,t$ & \unit{\milli\joule\per\centi\metre\squared} & 1,56 & 3,11 & 4,67 & 6,22 & 7,78 & 9,33 & 10,89 & 15,55 \\
    Épaisseur restante $e$ & \unit{\micro\metre} & 0,008 & 0,007 & 0,003 & 0,005 & 0,013 & 0,006 & 0,007 & 0,006 \\
    Épaisseur normalisée $e/\ez$ & --- & 0,006 & 0,005 & 0,002 & 0,004 & 0,010 & 0,005 & 0,006 & 0,004 \\
    \bottomrule
  \end{tabularx}
\end{table}

% --- Style commun des graphes ----------------------------------------
\pgfplotsset{
  topic1/.style={
    width=0.5\linewidth, height=3.6cm,
    tick label style={font=\scriptsize}, label style={font=\footnotesize},
    title style={font=\footnotesize, yshift=-3pt},
    /pgf/number format/use comma,
    grid=major, grid style={gray!20},
    tick align=outside, tick pos=left,
  },
  contraste/.style={
    topic1, xmin=1, xmax=30, ymin=0, ymax=1.05,
    xtick={1,2,5,10,20,30}, xticklabels={1,2,5,10,20,30},
    log ticks with fixed point,
    ytick={0,0.2,...,1}, yticklabel style={/pgf/number format/fixed, /pgf/number format/precision=1},
  },
}

\begin{figure}[H]
  \centering
  \begin{tikzpicture}
  \begin{groupplot}[
    group style={group size=2 by 2, horizontal sep=1.2cm, vertical sep=1.25cm,
                 xlabels at=edge bottom, ylabels at=edge left},
    contraste, xmode=log,
    xlabel={Dose $D$ (\unit{\milli\joule\per\centi\metre\squared})},
    ylabel={$e/\ez$},
  ]
    \nextgroupplot[title={Wafer 1 (\qty{2000}{\tour\per\minute}, $\ez = \qty{1,82}{\micro\metre}$)}]
    \addplot[gray, dashed, thin] coordinates {(3.217,1) (19.157,0)};
    \draw[gray, thin] (axis cs:19.157,0) -- (axis cs:19.157,0.28);
    \node[anchor=north east, font=\scriptsize, align=right] at (rel axis cs:0.97,0.95) {$\Dz \approx \qty{19}{\milli\joule\per\centi\metre\squared}$\\ $\gamma \approx 1{,}3$};
    \addplot[black, mark=*, mark size=1.3pt, thin] coordinates {(6.2200,0.6187) (9.3300,0.4277) (12.4400,0.2416) (15.5500,0.0950) (18.6600,0.0242) (21.7700,0.0029) (24.8800,0.0021) (27.9900,0.0009)};
    \nextgroupplot[title={Wafer 2 (\qty{3000}{\tour\per\minute}, $\ez = \qty{1,42}{\micro\metre}$)}]
    \addplot[gray, dashed, thin] coordinates {(2.866,1) (13.451,0)};
    \draw[gray, thin] (axis cs:13.451,0) -- (axis cs:13.451,0.28);
    \node[anchor=north east, font=\scriptsize, align=right] at (rel axis cs:0.97,0.95) {$\Dz \approx \qty{13}{\milli\joule\per\centi\metre\squared}$\\ $\gamma \approx 1{,}5$};
    \addplot[black, mark=*, mark size=1.3pt, thin] coordinates {(3.1100,0.9116) (6.2200,0.5674) (9.3300,0.2434) (12.4400,0.0109) (15.5500,0.0183) (18.6600,0.0077) (21.7700,0.0054) (24.8800,0.0043)};
    \nextgroupplot[title={Wafer 3 (\qty{4000}{\tour\per\minute}, $\ez = \qty{1,21}{\micro\metre}$)}]
    \fill[gray!18] (axis cs:1.555,0) rectangle (axis cs:3.110,1.05);
    \node[anchor=north west, font=\scriptsize] at (axis cs:1.555,1.02) {\Dz{} entre 1,6 et 3,1};
    \addplot[black, mark=*, mark size=1.3pt, thin] coordinates {(1.5550,0.5985) (3.1100,0.0233) (4.6650,0.0031) (6.2200,0.0053) (7.7750,0.0026) (9.3300,0.0000) (10.8850,0.0033) (15.5500,0.0049)};
    \nextgroupplot[title={Wafer 4 (\qty{5000}{\tour\per\minute}, $\ez = \qty{1,30}{\micro\metre}$)}]
    \fill[gray!18] (axis cs:1.000,0) rectangle (axis cs:1.555,1.05);
    \node[anchor=north west, font=\scriptsize] at (axis cs:1.000,1.02) {$\Dz < 1{,}6$};
    \addplot[black, mark=*, mark size=1.3pt, thin] coordinates {(1.5550,0.0064) (3.1100,0.0053) (4.6650,0.0024) (6.2200,0.0042) (7.7750,0.0098) (9.3300,0.0048) (10.8850,0.0056) (15.5500,0.0043)};
  \end{groupplot}
  \end{tikzpicture}
  \caption{Courbes de contraste de l'AZ1512 : épaisseur normalisée $e/\ez$ en fonction de la dose (échelle logarithmique). En pointillés, l'ajustement de la zone linéaire (wafers 1 et 2) ; en gris, l'encadrement de \Dz{} (wafers 3 et 4).}
  \label{fig:topic1-contraste}
\end{figure}

\begin{figure}[tb]
  \centering
  \begin{tikzpicture}
  \begin{groupplot}[
    group style={group size=3 by 1, horizontal sep=1.3cm},
    topic1, width=0.3\linewidth, height=4cm,
  ]
    % (a) épaisseur en fonction de la vitesse
    \nextgroupplot[title={(a) Courbe d'étalement},
      xlabel={Vitesse ($10^3$~\unit{\tour\per\minute})}, ylabel={\ez{} (\unit{\micro\metre})},
      xmin=1500, xmax=5500, ymin=1, ymax=2,
      xtick={2000,3000,4000,5000}, xticklabels={2,3,4,5},
      yticklabel style={/pgf/number format/fixed, /pgf/number format/precision=1}]
    \addplot[gray, dashed, thin, domain=1800:5200, samples=60] {78.66/sqrt(x)};
    \addplot[black, only marks, mark=*, mark size=1.3pt,
             error bars/.cd, y dir=both, y explicit]
      coordinates {(2000,1.81932) +- (0,0.03027) (3000,1.42422) +- (0,0.00397) (4000,1.21128) +- (0,0.07927) (5000,1.29558) +- (0,0.00296)};

    % (b) Dz en fonction de l'épaisseur initiale
    \nextgroupplot[title={(b) Sensibilité \Dz},
      xlabel={\ez{} (\unit{\micro\metre})}, ylabel={\Dz{} (\unit{\milli\joule\per\centi\metre\squared})},
      xmin=1.1, xmax=1.9, ymin=0, ymax=22,
      xtick={1.2,1.4,1.6,1.8},
      xticklabel style={/pgf/number format/fixed, /pgf/number format/precision=1}]
    \addplot[gray, dashed, thin, domain=1.1:1.9] {13.451+(x-1.42422)*14.442};
    \addplot[black, only marks, mark=*, mark size=1.3pt] coordinates {(1.81932,19.16) (1.42422,13.45)};
    \addplot[black, only marks, mark=*, mark size=1.3pt, error bars/.cd, y dir=both, y explicit]
      coordinates {(1.21128,2.333) +- (0,0.777)};
    \addplot[black, only marks, mark=triangle*, mark size=1.8pt, every mark/.append style={rotate=180}]
      coordinates {(1.29558,1.555)};
    \node[anchor=east, font=\scriptsize] at (axis cs:1.81932,19.16) {W1\,};
    \node[anchor=north west, font=\scriptsize] at (axis cs:1.42422,13.45) {W2};
    \node[anchor=south, font=\scriptsize] at (axis cs:1.21128,3.110) {W3};
    \node[anchor=west, font=\scriptsize] at (axis cs:1.29558,1.555) {W4};

    % (c) courbes attendues et mesurées des wafers 3 et 4
    \nextgroupplot[contraste, xmode=log, width=0.36\linewidth, height=4cm, title={(c) Wafers 3 et 4},
      xlabel={Dose $D$ (\unit{\milli\joule\per\centi\metre\squared})}, ylabel={$e/\ez$},
      xtick={1,3,10,30}, xticklabels={1,3,10,30}]
    \addplot[black, dashed, thin, domain=1:30, samples=120] {max(0,min(1,1.390*log10(10.376/x)))};
    \addplot[gray, dashed, thin, domain=1:30, samples=120] {max(0,min(1,1.390*log10(11.594/x)))};
    \addplot[black, mark=*, mark size=1.2pt, thin] coordinates {(1.5550,0.5985) (3.1100,0.0233) (4.6650,0.0031) (6.2200,0.0053) (7.7750,0.0026) (9.3300,0.0000) (10.8850,0.0033) (15.5500,0.0049)};
    \addplot[gray, mark=square*, mark size=1.1pt, thin] coordinates {(1.5550,0.0064) (3.1100,0.0053) (4.6650,0.0024) (6.2200,0.0042) (7.7750,0.0098) (9.3300,0.0048) (10.8850,0.0056) (15.5500,0.0043)};
  \end{groupplot}
  \end{tikzpicture}
  \caption{(a) Épaisseur initiale moyenne en fonction de la vitesse de rotation (barres : écart-type ; pointillés : loi en $1/\sqrt{\omega}$ ajustée sur les wafers 1 à 3). (b) Dose de clarification \Dz{} en fonction de l'épaisseur initiale ; pour les wafers 3 et 4, seul un encadrement est possible ; pointillés : tendance des wafers 1 et 2. (c) Courbes de contraste mesurées (traits pleins) et attendues d'après la tendance de \Dz{} et le $\gamma$ moyen des wafers 1 et 2 (pointillés), pour le wafer 3 (noir) et le wafer 4 (gris).}
  \label{fig:topic1-etalement-d100}
\end{figure}

\paragraph{Partie B.} \averifier{Angles de contact mesurés (tel que reçu, après déshydratation, après HMDS, Téflon) et observations au microscope du test d'adhésion, avec et sans HMDS.}

\subsection{Analyse}

\paragraph{Épaisseur et vitesse de rotation.}
Entre 2000 et \qty{4000}{\tour\per\minute}, l'épaisseur baisse de \qty{1,82}{\micro\metre} à \qty{1,21}{\micro\metre}, en suivant à peu près la loi attendue en $1/\sqrt{\omega}$ (figure~\ref{fig:topic1-etalement-d100}a). Le wafer~4 (\qty{1,30}{\micro\metre} à \qty{5000}{\tour\per\minute}) est plus épais que le wafer~3, ce qui n'est pas physique : nous pensons à une inversion de wafers ou à une vitesse mal réglée pendant le traitement en parallèle. Les films sont uniformes (écart-type de 3 à \qty{30}{\nano\metre}, tableau~\ref{tab:topic1-epaisseurs}), sauf le wafer~3 (\qty{79}{\nano\metre}) dont les zones O et A sont plus épaisses. Pour le wafer~1, l'écart-type vient surtout de la zone C, mesurée avec un GOF de 0,81 seulement.

\paragraph{Courbes de contraste des wafers 1 et 2.}
Les wafers 1 et 2 donnent des courbes complètes (figure~\ref{fig:topic1-contraste}). \Dz{} augmente avec l'épaisseur (\num{19} contre \qty{13}{\milli\joule\per\centi\metre\squared}, figure~\ref{fig:topic1-etalement-d100}b) : la résine absorbe les UV, donc le bas d'un film épais reçoit moins de dose. Le contraste reste proche ($\gamma = 1{,}3$ et $1{,}5$), car il dépend surtout de la résine et du développeur.

\paragraph{Courbes plates des wafers 3 et 4.}
C'est le point principal de la séance. En prolongeant la tendance des wafers 1 et 2, on attendait $\Dz \approx \qty{10}{\milli\joule\per\centi\metre\squared}$ pour le wafer~3 et $\approx \qty{12}{\milli\joule\per\centi\metre\squared}$ pour le wafer~4, avec $\gamma \approx 1{,}4$ : notre grille de doses (\qtyrange{1,6}{15,6}{\milli\joule\per\centi\metre\squared}) aurait alors couvert toute la transition (figure~\ref{fig:topic1-etalement-d100}c). Or la résine est déjà dissoute à \qty{3,1}{\milli\joule\per\centi\metre\squared} sur le wafer~3 et dès \qty{1,6}{\milli\joule\per\centi\metre\squared} sur le wafer~4 : on ne peut qu'encadrer \Dz, et $\gamma$ n'est pas mesurable. L'épaisseur seule n'explique pas un tel écart ; l'imprécision des temps très courts sur la lampe P.250 ou des différences de développement ont sans doute joué.

Le protocole donne des temps d'exposition à adapter selon les résultats des premiers wafers, mais nous avons traité les quatre wafers en parallèle avec l'assistante. Les wafers 3 et 4 ont donc reçu directement les temps suggérés (\qtyrange{0,5}{5}{\second}), sans correction possible, et presque toutes leurs zones ont dépassé \Dz. Il aurait fallu développer et mesurer chaque wafer avant d'exposer le suivant, puis recentrer les temps autour de la transition.

\subsection{Conclusion}
L'épaisseur de l'AZ1512 passe de \qty{1,82}{\micro\metre} à \qty{1,21}{\micro\metre} entre 2000 et \qty{4000}{\tour\per\minute}, conformément à la loi en $\omega^{-1/2}$ ; le wafer à \qty{5000}{\tour\per\minute} sort de cette tendance. Pour les deux films les plus épais, la dose de clarification augmente avec l'épaisseur (\qtylist{13;19}{\milli\joule\per\centi\metre\squared} pour \qtylist{1,42;1,82}{\micro\metre}) et le contraste vaut 1,3 à 1,5. Les films minces se sont dissous dès les plus faibles doses de notre grille, d'où des courbes plates : il faut traiter les wafers un par un et adapter les temps après une première zone test.